% ============================================================
%  SECTION III — DATASET CONSTRUCTION
%  \input'd from main.tex
% ============================================================

\section{Dataset Construction}
\label{sec:methodology}

This section describes the complete pipeline used to construct
\textbf{\datasetname} (\pathtt{Shanmuk4622/ai-detection-dataset-v2}).
The pipeline is implemented across eleven version-controlled Kaggle notebooks
(NB00--NB11), with a single shared library (\texttt{ai\_detect\_common.py},
pipeline version 1.2) that enforces a consistent schema and preprocessing
contract across all stages.
A single \texttt{MASTER\_SEED} value ($= 42$) controls all random state in
the build; changing it produces a statistically independent dataset with
the same structure.

%-------------------------------------------------------------------
\subsection{Real Image Collection}
\label{ssec:real}
%-------------------------------------------------------------------

The real-image corpus comprises 10,000 photographs drawn from two
complementary public datasets: 5,000 images from the MS-COCO 2017
validation split~\cite{lin2014coco} and 5,000 images from the
ImageNet-1k validation split~\cite{deng2009imagenet}.
The dual sourcing was deliberate: COCO images depict natural everyday
scenes with rich object co-occurrence and variable lighting, while
ImageNet provides a broader range of fine-grained categories across
all 1,000 synsets.
Together they prevent the real-image distribution from being dominated
by any single scene type, which would allow a detector to exploit
scene statistics rather than synthesis artefacts.

Images were downloaded from the canonical HuggingFace dataset
repositories (\texttt{detection-datasets/coco} and
\texttt{ILSVRC/imagenet-1k}) and subjected immediately to the canonical
preprocessing pipeline described in Section~\ref{ssec:preprocess},
before being stored in the repository.
No post-hoc filtering based on image content was performed; only images
that failed lossless integrity checks (corrupt headers, truncated files)
were discarded, amounting to fewer than 0.1\% of the downloaded corpus.

%-------------------------------------------------------------------
\subsection{Caption Generation}
\label{ssec:captions}
%-------------------------------------------------------------------

A caption-grounded pairing strategy is the central mechanism by which
the dataset eliminates content as a confounding variable between real
and AI images.
Rather than generating prompts independently per generator — which
would introduce semantic distribution differences between classes — we
derive a single textual description from each real image and pass it
unchanged to all six generators.
This guarantees that the semantic content available to all generators
is held constant, so any remaining visual difference between real and
AI images must be attributable to the synthesis process itself.

We caption each real image using
\textbf{BLIP-2}~\cite{li2023blip2} with the
\pathtt{Salesforce/blip2-opt-2.7b} checkpoint, running beam search
with five beams, \texttt{max\_new\_tokens}~$= 60$, and a
\texttt{repetition\_penalty} of 1.3.
The raw caption is then cleaned by stripping a set of common
uninformative lead-in phrases (e.g., ``the image shows'', ``there is'',
``a photo of'') and restoring sentence case.
The cleaned caption is then token-counted using the
\pathtt{openai/clip-vit-large-patch14} tokeniser~\cite{radford2021clip}
and truncated to a maximum of 75 tokens.
This cap is set deliberately below the CLIP context limit of 77 tokens
to accommodate the special \texttt{[BOS]} and \texttt{[EOS]} tokens
added by each generator's text encoder.
The resulting prompt therefore fits within the text encoder budget of
all six generators without padding or error.

The full verifier in NB02 confirmed that, after captioning:
all 10,000 real images were assigned exactly one caption,
zero captions were empty or duplicated, and
zero captions exceeded 77 CLIP tokens.
These captions are stored in the \texttt{captions/} prefix of the
dataset repository as Apache Parquet files and are publicly accessible
alongside the image data.

%-------------------------------------------------------------------
\subsection{Synthetic Image Generation}
\label{ssec:generation}
%-------------------------------------------------------------------

\subsubsection{Generator Selection}

The six generators were selected to span the major architectural
families and training paradigms active in the text-to-image literature
as of mid-2024.
Table~\ref{tab:generators} lists each generator with its model
identifier, backbone architecture, native resolution, and the inference
hyperparameters used during generation.

\begin{table*}[t]
\centering
\caption{Generator configurations used to produce the synthetic
portion of \textbf{\datasetname}.
All images are generated at the listed native resolution and then
downsampled to $512 \times 512$ by the canonical preprocessing
pipeline. Guidance~$= 0.0$ for FLUX.1-schnell reflects its
distilled, guidance-free inference regime. W\"{u}rstchen uses a
two-stage prior--decoder scheme with separate step and guidance
parameters for each stage.}
\label{tab:generators}
\small
\setlength{\tabcolsep}{5pt}
\begin{tabular}{@{}l>{\raggedright\arraybackslash}p{3.5cm}>{\raggedright\arraybackslash}p{3.6cm}cccl@{}}
\toprule
Key & Model & Architecture & Native (px) & Steps & Guidance & Ref. \\
\midrule
\texttt{sd15}
  & stable-diffusion-v1-5
  & UNet LDM
  & 512
  & 20
  & 7.0
  & \cite{rombach2022ldm} \\
\texttt{sdxl}
  & stable-diffusion-xl-base-1.0
  & UNet LDM (3$\times$ larger)
  & 1024
  & 8
  & 7.0
  & \cite{podell2024sdxl} \\
\texttt{flux\_schnell}
  & FLUX.1-schnell
  & Rectified flow transformer
  & 1024
  & 4
  & 0.0
  & \cite{flux2024bfl} \\
\texttt{kandinsky22}
  & kandinsky-2-2-decoder
  & Prior + UNet LDM
  & 768
  & 25
  & 4.0
  & \cite{razzhigaev2023kandinsky} \\
\texttt{pixart\_sigma}
  & PixArt-Sigma-XL-2-1024-MS
  & Diffusion Transformer (DiT)
  & 1024
  & 20
  & 4.5
  & \cite{chen2024pixart} \\
\texttt{wuerstchen}
  & warp-ai/wuerstchen
  & Hierarchical LDM (42$\times$ compression)
  & 1024
  & 20 / 10
  & 4.0 / 0.0
  & \cite{pernias2023wuerstchen} \\
\bottomrule
\end{tabular}
\end{table*}

The selection covers three distinct backbone paradigms: (i) \emph{UNet
latent diffusion models}~\cite{rombach2022ldm}, represented by SD~1.5,
SDXL, and the decoder stage of Kandinsky~2.2; (ii) \emph{attention-based
diffusion transformers}~\cite{chen2024pixart}, represented by
PixArt-$\Sigma$ and FLUX.1-schnell; and (iii) a \emph{hierarchical
two-stage architecture}, represented by W\"{u}rstchen.
This diversity is intentional: prior work has shown that forensic
artefacts are shaped primarily by model architecture rather than by
fine-tuning or personalisation~\cite{hong2024wildfake}, and covering all
three families ensures that the dataset supports evaluation of detectors
beyond the UNet-LDM regime that dominates most existing benchmarks.

\subsubsection{Deterministic Seed Assignment}

Each (generator, real image) pair is assigned a deterministic inference
seed computed as:
\begin{equation}
\begin{split}
s(m, r) = \big(\mathrm{int}(\mathrm{SHA\text{-}256}(
          &\,\texttt{SEED} : m : r)_{[0:15]},\, 16)\big) \\
          &\bmod\, (2^{31} - 1)
\end{split}
\label{eq:seed}
\end{equation}
where \texttt{SEED} denotes \texttt{MASTER\_SEED}, $m$ is the generator
key, $r$ is the \texttt{source\_real\_id}, the colon denotes string
concatenation, and the first 15 hexadecimal digits of the SHA-256
digest are interpreted as an integer.
This produces a unique, reproducible seed for every image without any
risk of collision between generators.
Changing \texttt{MASTER\_SEED} in NB00 regenerates the entire dataset
with a statistically independent set of images while preserving the
pipeline structure.

\subsubsection{Inference Settings}

All models were loaded with \texttt{torch.float16} on a single NVIDIA
Tesla T4 GPU (15.6 GB VRAM) using the HuggingFace Diffusers library,
with \texttt{safety\_checker=None} and \texttt{add\_watermarker=False}
to avoid any post-processing that could introduce systematic
class-specific artefacts.
The \texttt{add\_watermarker=False} flag is particularly important:
invisible watermarking schemes applied at the pipeline level (e.g.,
the default invisible watermark in SDXL) would introduce a machine-readable
signal that a detector could exploit without learning any synthesis-related
visual feature.

SDXL was run with 8 DPM++ steps to reflect real-world fast-inference
usage rather than the slower default schedule.
FLUX.1-schnell used 4 steps with classifier-free guidance scale $= 0.0$,
consistent with its distilled, guidance-free inference regime.
W\"{u}rstchen used a two-stage procedure: the prior model
(\pathtt{stabilityai/warp-ai/wuerstchen-prior}) ran for 20 steps with
guidance~$= 4.0$ to produce image embeddings, followed by the decoder
(\pathtt{stabilityai/warp-ai/wuerstchen}) for 10 steps with
guidance~$= 0.0$.
All other generators used standard scheduler settings from their
respective official repositories.

For each real image, the pipeline generated exactly one synthetic image
per generator, producing $10{,}000 \times 6 = 60{,}000$ AI images in total.

%-------------------------------------------------------------------
\subsection{Canonical Preprocessing Pipeline}
\label{ssec:preprocess}
%-------------------------------------------------------------------

The canonical preprocessing pipeline is the central technical
contribution of this work.
It is implemented as a single Python function
(\texttt{canonical\_preprocess}) in the shared library and applied
identically to every image — real or AI — with no branching on the
class label.
The invariance to label is enforced both by code structure
(the function takes only a PIL image as input) and by automated
testing (the verifier in NB10 checks that all stored images pass
the \texttt{png\_is\_canonical} predicate before the split is
finalised).

The pipeline consists of the following steps, applied in order:

\begin{enumerate}

\item \textbf{EXIF transpose.}
Any rotation metadata embedded in the JPEG header is applied and the
EXIF block is discarded.
This prevents AI images — which carry no EXIF data — from being
trivially distinguished from real photographs by the presence of
orientation tags.

\item \textbf{RGB conversion.}
The image is converted to three-channel RGB, normalising palette,
greyscale, and RGBA inputs to a common colour space.

\item \textbf{Centre-crop to square.}
The shorter side of the image defines the crop dimension.
Formally, for an image of width $W$ and height $H$ with
$S = \min(W, H)$, the crop region is
$\left[\frac{W-S}{2},\, \frac{H-S}{2},\,
 \frac{W+S}{2},\, \frac{H+S}{2}\right]$.
This avoids aspect-ratio padding, which would introduce blank regions
distinguishable by texture statistics.

\item \textbf{Lanczos resize to $512 \times 512$.}
The square crop is downsampled using Lanczos interpolation
(\texttt{Image.LANCZOS}).
Lanczos is chosen because it is the standard high-quality downsampling
filter; using different resampling methods for real and AI images would
produce class-specific aliasing patterns.

\item \textbf{JPEG equalisation.}
The resized image is re-encoded to JPEG at quality factor 95
with 4:4:4 chroma subsampling (\texttt{subsampling=0}) and
immediately decoded back to a pixel array.
This step is the primary defence against the compression bias
identified by Grommelt~\etal~\cite{grommelt2024jpeg}: real images
from COCO and ImageNet carry varying JPEG quality factors (modal
$Q \approx 96$ in ImageNet), while generated images are natively
lossless.
By passing every image through the same JPEG codec at the same quality
factor, we equalise the compression history so that no detector can
distinguish real from AI by measuring DCT coefficient distributions.
The round-trip through a fixed quality factor also eliminates
ICC colour profiles and other ancillary metadata.

\item \textbf{PNG serialisation.}
The equalised pixel array is written to a lossless PNG with
\texttt{compress\_level=6} and no EXIF, ICC, or chunk metadata.
PNG is chosen as the storage format because (a) it is lossless, so
no further compression artefacts are introduced during retrieval, and
(b) the absence of metadata eliminates any remaining header-based
shortcut.

\end{enumerate}

Every stored image is verified against the predicate
\texttt{png\_is\_canonical}, which checks that the format is PNG,
the mode is RGB, the spatial dimensions are $512 \times 512$, and
no ICC profile is present.
A SHA-256 digest of the PNG bytes is stored alongside each image in the
schema and re-verified during the leak audit (NB10) to detect any
post-hoc modification of stored images.

The full preprocessing algorithm is reproduced in
Algorithm~\ref{alg:preprocess} for clarity and reproducibility.

\begin{algorithm}[t]
\caption{Canonical preprocessing (applied identically to real and AI images)}
\label{alg:preprocess}
\begin{algorithmic}[1]
\Require PIL image $I$, target size $T = 512$, JPEG quality $Q = 95$
\State $I \leftarrow \textsc{ExifTranspose}(I)$
\State $I \leftarrow \textsc{ConvertRGB}(I)$
\State $W, H \leftarrow \textsc{Size}(I)$;
       $S \leftarrow \min(W, H)$
\State $I \leftarrow \textsc{Crop}\!\left(I,\;\tfrac{W-S}{2},\;\tfrac{H-S}{2},\;S,\;S\right)$
\State $I \leftarrow \textsc{Resize}(I, T \times T, \text{Lanczos})$
\State $I \leftarrow \textsc{JpegDecode}(\textsc{JpegEncode}(I, Q, \text{ss}{=}0))$
  \Comment{JPEG equalisation, 4:4:4}
\State \Return $\textsc{PngEncode}(I, \text{compress\_level}{=}6)$
\end{algorithmic}
\end{algorithm}

%-------------------------------------------------------------------
\subsection{Dataset Schema}
\label{ssec:schema}
%-------------------------------------------------------------------

All images and metadata are stored as Apache Parquet files using the
ZSTD compression codec.
Table~\ref{tab:schema} lists the full schema.
The \texttt{source\_real\_id} field is the dataset's pairing key:
for real images it equals \texttt{image\_id}, and for every AI image
it points back to the real image whose caption was used as the
generation prompt.
Researchers can therefore reconstruct any real--AI pair from a single
join on this field.

Two provenance fields, \texttt{orig\_width} and \texttt{orig\_height},
record the pre-crop pixel dimensions of each image.
These fields are retained strictly for provenance; they encode
information about the original source and must not be used as training
features, since doing so would reintroduce a resolution-based shortcut.

\begin{table}[t]
\centering
\caption{Dataset schema. Fields marked $\dagger$ must not be used as
training features (they encode source-level provenance).}
\label{tab:schema}
\footnotesize
\setlength{\tabcolsep}{4pt}
\begin{tabular}{@{}>{\raggedright\arraybackslash}p{2.5cm}l>{\raggedright\arraybackslash}p{3.4cm}@{}}
\toprule
Column & Type & Description \\
\midrule
\texttt{image\_id} & \texttt{string} & Globally unique identifier \\
\texttt{source\_real\_id} & \texttt{string} & Pairing key; $=$ \texttt{image\_id} for real rows \\
\texttt{label} & \texttt{int8} & 0 = real, 1 = AI \\
\texttt{generator} & \texttt{string} & Generator key (Table~\ref{tab:generators}) \\
\texttt{source\_dataset} & \texttt{string} & \texttt{coco}, \texttt{imagenet}, or generator key \\
\texttt{prompt} & \texttt{string} & BLIP-2 caption; \texttt{null} for real rows \\
\texttt{image} & \texttt{binary} & Canonical $512\!\times\!512$ RGB PNG bytes \\
\texttt{width}, \texttt{height} & \texttt{int16} & Always 512 \\
\texttt{orig\_width}$^\dagger$ & \texttt{int32} & Pre-crop width \\
\texttt{orig\_height}$^\dagger$ & \texttt{int32} & Pre-crop height \\
\texttt{gen\_model\_id} & \texttt{string} & HuggingFace model identifier \\
\texttt{gen\_steps} & \texttt{int16} & Inference steps used \\
\texttt{gen\_guidance} & \texttt{float32} & Classifier-free guidance scale \\
\texttt{gen\_native\_res} & \texttt{int16} & Native generation resolution \\
\texttt{seed} & \texttt{int64} & Per-image deterministic seed (Eq.~\ref{eq:seed}) \\
\texttt{caption\_model} & \texttt{string} & Captioner checkpoint \\
\texttt{pipeline\_version} & \texttt{string} & Library version (\texttt{1.2}) \\
\texttt{sha256} & \texttt{string} & Hex SHA-256 of PNG bytes \\
\texttt{created\_utc} & \texttt{string} & ISO-8601 UTC creation timestamp \\
\bottomrule
\end{tabular}
\end{table}

%-------------------------------------------------------------------
\subsection{Dataset Split}
\label{ssec:splits}
%-------------------------------------------------------------------

Splits are assigned at the \emph{pair level}: a \texttt{source\_real\_id}
and all seven images associated with it (one real, six AI) are
assigned to the same fold.
This prevents any possibility of a real image appearing in training
while its AI partners appear in the test set, which would leak prompt
content and scene composition across the split boundary.

The split assignment function maps each \texttt{source\_real\_id} to a
uniformly distributed scalar in $[0, 1)$ via:
\begin{equation}
f(r) = \frac{\mathrm{int}(\mathrm{SHA\text{-}256}
        (\texttt{SEED} : r)_{[0:8]},\, 16)}{100{,}000}
\bmod 1
\label{eq:split}
\end{equation}
Using the assignment of Equation~\ref{eq:split}, pairs with
$f(r) < 0.70$ are assigned to \emph{train}, those with
$0.70 \leq f(r) < 0.85$ to \emph{validation}, and the remainder to
\emph{test}.
The resulting image counts — 49,392 train, 10,122 validation,
10,486 test — are within 0.5\% of the target $70/15/15$ allocation,
confirming that the hash function is well-calibrated to the
\texttt{MASTER\_SEED}.
Table~\ref{tab:split_counts} shows the per-generator breakdown.

\begin{table}[t]
\centering
\caption{Image counts per split and generator.
All generators contribute exactly 10,000 images.
The real class contributes 10,000 images total.}
\label{tab:split_counts}
\small
\begin{tabular}{lrrr|r}
\toprule
Generator & Train & Val & Test & Total \\
\midrule
Real             & 7{,}056 & 1{,}446 & 1{,}498 & 10{,}000 \\
SD~1.5           & 7{,}056 & 1{,}446 & 1{,}498 & 10{,}000 \\
SDXL             & 7{,}056 & 1{,}446 & 1{,}498 & 10{,}000 \\
FLUX.1-schnell   & 7{,}056 & 1{,}446 & 1{,}498 & 10{,}000 \\
Kandinsky~2.2    & 7{,}056 & 1{,}446 & 1{,}498 & 10{,}000 \\
PixArt-$\Sigma$  & 7{,}056 & 1{,}446 & 1{,}498 & 10{,}000 \\
W\"{u}rstchen       & 7{,}056 & 1{,}446 & 1{,}498 & 10{,}000 \\
\midrule
\textbf{Total}   & \textbf{49,392} & \textbf{10,122} & \textbf{10,486} & \textbf{70,000} \\
\bottomrule
\end{tabular}
\end{table}

Split assignments are materialised in \texttt{splits.parquet}
(10,000 rows, one per real image) and replicated in
\texttt{manifest.parquet} (70,000 rows, one per image across all
generators).
The manifest includes \texttt{image\_id}, \texttt{source\_real\_id},
\texttt{generator}, \texttt{label}, \texttt{sha256}, and \texttt{split}
for every image, enabling efficient split-stratified data loading
without reading any image bytes.

%-------------------------------------------------------------------
\subsection{Deduplication}
\label{ssec:dedup}
%-------------------------------------------------------------------

After all generation notebooks completed, NB09 aggregated metadata
across all 70,000 images and checked for duplicate \texttt{image\_id}s
and for groups of images that share identical SHA-256 digests across
different class labels.
No duplicate \texttt{image\_id}s were found.
One group of four W\"{u}rstchen images shared identical byte content
(the same SHA-256 digest), a consequence of the generator collapsing
to a degenerate output for those particular prompts.
These images retain distinct \texttt{image\_id}s in the schema and
are flagged in \texttt{manifest.parquet} via their shared
\texttt{sha256} value so that users can exclude them if desired.
No cross-label byte collisions (a real image identical in bytes to an
AI image) were found.

%-------------------------------------------------------------------
\subsection{Summary Statistics}
\label{ssec:summary}
%-------------------------------------------------------------------

Table~\ref{tab:summary} summarises the final dataset.
The complete build pipeline spans eleven notebooks, executing
sequentially from NB00 to NB11, with generation notebooks NB03--NB08
parallelised across multiple Kaggle free-tier accounts.
Parallelisation uses the hash-based worker-partitioning scheme
described in Section~\ref{ssec:generation}: each Kaggle account
processes a shard of real images assigned by
$\mathrm{SHA}\text{-}256(\texttt{sid}) \bmod N_{\text{workers}}$,
where \texttt{sid} is the account identifier and $N_{\text{workers}}$
is the number of parallel accounts.
This scheme guarantees that each (real image, generator) pair is
processed by exactly one account, preventing duplicate generation
without requiring inter-process coordination.
The total generation time was approximately 180~GPU-hours across
single NVIDIA Tesla~T4 instances.

\begin{table}[t]
\centering
\caption{Summary statistics for \textbf{\datasetname}.}
\label{tab:summary}
\setlength{\tabcolsep}{4pt}
\begin{tabular}{@{}l>{\raggedright\arraybackslash}p{5.15cm}@{}}
\toprule
Property & Value \\
\midrule
Total images                  & 70,000 \\
Real images                   & 10,000 (5,000 COCO + 5,000 ImageNet) \\
AI images                     & 60,000 (10,000 per generator) \\
Generators                    & 6 \\
Resolution                    & $512 \times 512$ RGB PNG \\
Caption model                 & BLIP-2 (\texttt{blip2-opt-2.7b})~\cite{li2023blip2} \\
Max prompt tokens             & 75 (CLIP tokeniser)~\cite{radford2021clip} \\
Pipeline version              & 1.2 \\
Master seed                   & 42 \\
Split ratio                   & 70 / 15 / 15 (pair-level) \\
Train / Val / Test            & 49,392 / 10,122 / 10,486 \\
Storage format                & Apache Parquet (ZSTD) \\
Repository                    & \texttt{Shanmuk4622/}\allowbreak\texttt{ai-detection-dataset-v2} \\
Licence                       & Non-commercial research (most restrictive applicable) \\
\bottomrule
\end{tabular}
\end{table}


% ============================================================
% ADDITIONAL BIBLIOGRAPHY ENTRIES (add to your .bib / thebibliography)
% ============================================================

% ── Add these entries to the thebibliography block already created
%    in related_work_and_references.tex ────────────────────────────
